\chapter{Tests of the Standard Model}
\label{ch:tests-of-sm}

In combination with experimental measurements, lattice QCD is used extensively to test the Standard Model.
This is perhaps the most successful application of LQFT and the broad program that is undertaken by the community is summarised in the frequent reports of the  \textit{Flavour Lattice Averaging Group (FLAG)} to which the reader is referred for latest averages of particular quantities and historical evolution of these results.

In this chapter, we present examples of this  program that highlight particular themes that are broadly  applicable.

\todo{The parameters of QCD themselves --- $\alpha_s$ and the quark masses ---
are determinations rather than tests, and are now in
Ch.~\ref{ch:strong-coupling-sf}. This chapter is only about quantities that,
combined with experiment, over-determine something.}

\section{The CKM matrix and test of unitarity}

\subsection{First-row unitarity}

\todo{New stub: $V_{ud}$ and $V_{us}$ from $K_{\ell 2}$ and $K_{\ell 3}$
against superallowed nuclear $\beta$ decay; the lattice inputs are
$f_K/f_\pi$ and $f_+^{K\pi}(0)$. This is the cleanest example in the chapter of
a closed test --- three numbers that must sum to one --- and it is worth
setting up as the template before the triangle.}

\subsection{The unitarity triangle}

\todo{New stub: how each lattice quantity enters, and where the constraints
disagree. Distinguish the angles (experimental) from the sides (lattice).}

\section{Neutral meson mixing}

\todo{New section: $B_K$ and $\varepsilon_K$; $\Delta M_d$, $\Delta M_s$ and the
bag parameters with the $SU(3)$-breaking ratio $\xi$; $D$-mixing and why it is
harder. This is a distinct constraint from the CKM elements of the previous
section --- box diagrams rather than tree-level decays --- and it is the place
where four-quark operator renormalisation and mixing does its real work
(cross-reference Ch.~\ref{ch:matrix-elements}).}

\section{CP violation and kaon decays}

\subsection{The $\Delta I=1/2$ rule and its non-perturbative origin}

\subsection{$\epsilon'/\epsilon$}

\section{Other kaon physics}

\todo{New stub, as a home for the kaon quantities that are neither $K\to\pi\pi$
nor CKM inputs. Candidates: $K\to\pi\nu\bar\nu$; the $K_L$--$K_S$ mass
difference; long-distance contributions generally, which is what makes these
hard --- they need bi-local operators and a genuinely different technique from
anything else in the chapter, and that is the transferable theme. Rare and
radiative kaon decays if there is room.}

\section{Decays of heavy hadrons}

\todo{Keep distinct from the CKM section above: that one is organised around
unitarity, this one around the decays themselves. Semileptonic form factors and
decay constants; the inclusive--exclusive tension in $|V_{cb}|$ and $|V_{ub}|$
and what the lattice can now say about the inclusive side (which is an inverse
problem --- cross-reference Ch.~\ref{ch:data-analysis},
\S\ref{sec:inverse-problems}); lepton flavour universality, $R(D)$ and
$R(D^{(*)})$; heavy-quark discretisations as the common systematic.}

\subsection{FCNC and rare decays}

\todo{New stub: $B\to K^{(*)}\ell\ell$ form factors, $B_s\to\mu\mu$, and the
$b\to s\ell\ell$ observables.}

\section{The anomalous magnetic moment of the muon}

\todo{Deliberately kept to one section. The hadronic vacuum polarisation and
light-by-light contributions, window observables and why they were introduced,
and the role of isospin breaking (Ch.~\ref{ch:isospin-qed}) --- as an
illustration of what sub-percent precision demands, not as a survey of a moving
target. Refer to the current Theory Initiative and FLAG reports for status.}

\section{Fundamental symmetry tests}

\subsection{Neutron EDM}

\subsection{Neutrinoless double-$\beta$ decay}

\subsection{Matrix elements for direct dark matter detection}

\todo{New stub: the nucleon $\sigma$-terms and the strange and charm content of
the nucleon, which set the scattering cross sections that direct-detection
limits are quoted against. A good example of lattice input to a search with no
accelerator attached.}

\subsection{Nucleon charges and precision $\beta$ decay}

\todo{New stub: $g_A$ as the field's benchmark quantity and a useful cautionary
history; the scalar and tensor charges $g_S$, $g_T$ as the inputs that convert
neutron-decay measurements into limits on new interactions.}

\subsection{Proton decay}

\todo{New stub: matrix elements of the baryon-number-violating three-quark
operators, which turn a GUT into a lifetime prediction.}

\section{How averages are made}

\todo{New stub, and optional --- delete if the chapter is already long enough.
The FLAG quality criteria, what makes results correlated, and what a
``tension'' means when the inputs share ensembles. The chapter opens by
pointing the reader at FLAG; explaining the machinery rather than only citing
it would be unusual for a textbook and is a genuinely transferable skill.}
